\section{Conclusion}
This draft presented \tempo{} as a lightweight synchronous runtime and explained
how its control constructs are implemented with OCaml 5.3 algebraic effects.
The report positioned \tempo{} within the reactive control lineage, clarified the
core constructs (logical instants, signals, guards, weak preemption), and
detailed the scheduler and effect-handler design. The current level of
confidence in correctness relies on an extensive test suite that exercises the
expected instant-based semantics across all public features (signals, aggregate
signals, guards, preemption, await\_immediate, parallel composition, and
derived combinators).
See Section~\ref{sec:feature-tests} for a feature-to-test mapping.

The determinism guarantees discussed in this report assume standard purity
discipline: programs should avoid shared mutable state whose outcomes depend on
intra-instant scheduling, and aggregate combine functions should be
order-insensitive. These assumptions match the synchronous control model and
clarify the scope of the current guarantees.

\paragraph{Future work}\label{sec:future-work}
We plan to formalize correctness in two steps. First, we will define a
high-level synchronous semantics and prove a refinement relation with an
abstract scheduler model (e.g., along the lines of formal treatments of
synchronous languages and their operational models
\cite{esterel1985,esterel1992,reactiveml2008}). Second, we will connect this
abstract scheduler to the concrete effect-handler implementation, drawing on
formal work on algebraic effects and handlers
\cite{plotkinpretnar2009,eff2013}. This two-layer approach is intended to
bridge the gap between the user-facing semantics and the low-level execution
mechanism.

Two additional concerns are central. First, the use of references for
communication can break semantic determinism: even if the implementation fixes
the evaluation order within an instant, that order is not specified by the
semantics and can change across versions. Addressing this requires static
analyses that detect or rule out order-sensitive behaviors. Second, \tempo{}'s
parallelism is currently logical rather than physical. Supporting physical
parallelism risks reintroducing nondeterminism unless the scheduler enforces a
deterministic memory discipline. We plan to investigate recent work by Pouzet
and Mendler on PSM for guidance on deterministic shared-memory parallelism
\cite{mendlerpouzet2025psm} (currently a preprint). Finally, adding asynchronous
computations for selected programming idioms is an open direction, provided
their interaction with instant-based determinism remains controlled.

Additional avenues include: (1) formalizing the host I/O contract with a
precise notion of input injection and output observation per instant, aligned
with synchronous language interfaces \cite{lustre1991,esterel1992}; (2)
designing static analyses or effect disciplines to detect order-sensitive use
of shared mutable state, in the spirit of clock/causality analyses
\cite{lustre1991,esterel1992}; (3) exploring partial compilation or macro
expansion to reduce effect overhead in hot paths, following compilation
strategies in synchronous languages \cite{esterel1992}; (4) adding instant-level
profiling support (queue sizes, wakeups, guard checks) to aid performance
diagnosis, building on tooling experience from reactive implementations
\cite{reactiveml2015}; and (5) clarifying aggregation semantics for
non-commutative combines, either by fixing a reduction order or by restricting
such combines to controlled modes, drawing on deterministic stream models
\cite{kahn1974}.

As a simple illustration, two parallel tasks that both mutate a shared
reference and then emit a signal can produce different outcomes depending on
their evaluation order within the instant. Since that order is not part of the
semantic contract, relying on it compromises determinism and motivates static
analyses that forbid or control such patterns.
\begin{lstlisting}[style=tempo]
let r = ref 0 in
parallel [
  (fun () -> r := 1; emit s !r);
  (fun () -> r := 2; emit s !r)
]
\end{lstlisting}
Depending on the evaluation order in the instant, the emitted value can be
\texttt{1} or \texttt{2}, despite identical source code.
